\chapter{Konstrukcja Pontriagina--Thoma i podstawy teorii kobordyzmów}
\label{ch:pt-v9}

\section{Motywacja: jak odwzorowanie w sferę koduje podrozmaitość?}
W poprzednim rozdziale zdefiniowaliśmy stabilne grupy homotopii sfer
za pomocą zawieszeń. Ich elementy mają interpretację geometryczną:
można je reprezentować przez zwarte podrozmaitości o \emph{uramowanej}
wiązce normalnej. Przejście od podrozmaitości do odwzorowania
polega na ściągnięciu do jednego punktu wszystkiego poza jej
otoczeniem tubularnym. Jest to \emph{konstrukcja Pontriagina--Thoma}.
Po zastąpieniu uramowania dowolną wiązką normalną otrzymujemy
przestrzenie Thoma i widma $MO$, $MSO$; służą one do opisu
kobordyzmów nieorientowanych i zorientowanych. Odróżniamy przy tym
trzy poziomy: twierdzenie dla uramowanych podrozmaitości,
jego wersję dla wiązek normalnych i twierdzenia obliczeniowe
o pierścieniach kobordyzmu.

Pracujemy w kategorii \emph{gładkich}, Hausdorffowskich,
drugoprzeliczalnych rozmaitości. Rozmaitość \emph{zamknięta}
oznacza tutaj zwartą rozmaitość bez brzegu; dopuszczamy
rozmaitość pustą. Podstawy o podrozmaitościach, transwersalności
i otoczeniach tubularnych przypomniano w dodatku. Rozkłady
kobordyzmów na uchwyty są przedmiotem osobnego skryptu z geometrii
różniczkowej; tutaj badamy przede wszystkim \emph{niezmienniki}
kobordyzmu.

\section{Włókno wartości regularnej i uramowanie normalne}
\begin{definicja}[Wiązka normalna i uramowanie]
Niech $M^n\subseteq Q^{n+k}$ będzie domkniętą, gładko
osadzoną podrozmaitością, $k\ge1$. Jej \emph{wiązką normalną}
nazywamy rzeczywistą wiązkę wektorową rzędu $k$,
$\nu_{M/Q}:=TQ|_M/TM$. Po wyborze metryki można ją
utożsamić z prostopadłym dopełnieniem $TM$ w $TQ|_M$.
\emph{Uramowaniem normalnym} nazywamy izomorfizm wiązek
$\varphi:\nu_{M/Q}\to M\times\R^k$, czyli ciągłą
(gładką) rodzinę uporządkowanych baz we włóknach.
\emph{Podrozmaitość uramowana} oznacza parę $(M,\varphi)$.
Zmiana uramowania nie musi zachowywać jego klasy kobordyzmu.
\end{definicja}

\begin{lemat}[Wartość regularna i uramowanie]
\label{lem:pt-regular-v9}
Niech $Q^{n+k}$ będzie zwartą gładką rozmaitością,
$f:Q\to S^k$ odwzorowaniem gładkim, a $q\in S^k$
wartością regularną $f$. Wtedy $M=f^{-1}(q)$
jest zamkniętą podrozmaitością wymiaru $n$.
Po wyborze uporządkowanej bazy przestrzeni $T_qS^k$
różniczka $f$ wyznacza kanoniczne \emph{względem tego
wyboru} uramowanie $\nu_{M/Q}$.
\end{lemat}
\begin{proof}
Dla $x\in M$ surjektywność
$\D f_x:T_xQ\to T_qS^k$ daje, z twierdzenia o wartości
regularnej, gładką podrozmaitość $M$ z
$T_xM=\ker\D f_x$. Zatem $\D f_x$ indukuje
izomorfizm ilorazów
$T_xQ/T_xM\to T_qS^k$; izomorfizmy zależą gładko od
$x$ i składają się na izomorfizm wiązki
$\nu_{M/Q}\cong M\times T_qS^k$.
Ustalona baza $T_qS^k\cong\R^k$ kończy konstrukcję.
Zwartość $Q$ zapewnia zwartość zbioru $M$, a jego
domkniętość wynika z ciągłości $f$.
\end{proof}

\begin{uwaga}[Przedmiot uramowania]
Uramowujemy \emph{wiązkę normalną}, niekoniecznie
wiązkę styczną. Istnienie stabilnego uramowania normalnego
jest równoważne istnieniu stabilnego uramowania stycznego,
lecz konkretne uramowania należy porównać za pomocą relacji
$TM\oplus\nu_{M/\R^{n+k}}\cong\underline\R^{n+k}$.
W przestrzeni $S^{n+k}$ używamy wykresu
$S^{n+k}\setminus\{\infty\}\cong\R^{n+k}$;
w ten sposób ramowanie normalne nie zależy od trywialności
całej wiązki stycznej sfery.
\end{uwaga}

\section{Przestrzeń Thoma i odwzorowanie ściągające}
\begin{definicja}[Przestrzeń Thoma]
Niech $\xi\to B$ będzie rzeczywistą wiązką wektorową
rzędu $k$ z metryką włóknową. Oznaczmy przez $D(\xi)$
odpowiednio wiązkę dysków jednostkowych, a przez $S(\xi)$
wiązkę sfer jednostkowych. \emph{Przestrzenią Thoma} jest
przestrzeń punktowana
$\operatorname{Th}(\xi):=D(\xi)/S(\xi)$,
w której całą $S(\xi)$ utożsamiamy do jednego punktu.
Dla wiązki trywialnej
$\underline\R^k_B\to B$ mamy naturalny homeomorfizm
$\operatorname{Th}(\underline\R^k_B)\cong B_+\wedge S^k$;
$B_+$ oznacza $B$ z dodanym rozłącznym punktem bazowym.
\end{definicja}

\begin{lemat}[Odwzorowanie ściągające]
\label{lem:pt-collapse-v9}
Niech $i:M^n\hookrightarrow Q^{n+k}$ będzie gładkim
osadzeniem zwartej podrozmaitości we wnętrzu zwartej
rozmaitości $Q$ (ewentualnie z brzegiem), z
$M\cap\partial Q=\varnothing$. Wybierzmy otoczenie tubularne
$U\cong \operatorname{int}D(\nu_{M/Q})$ z domknięciem
we wnętrzu $Q$. Wtedy istnieje ciągłe punktowane
odwzorowanie
\[
c_{M,Q}: Q/\partial Q\longrightarrow
\operatorname{Th}(\nu_{M/Q}),
\]
przypisujące wszystkim punktom spoza $U$ punkt bazowy,
a punktom w $U$ ich współrzędne tubularne. Jeśli
$\partial Q=\varnothing$, przez $Q/\partial Q$ rozumiemy
$Q_+$ z oddzielnym punktem bazowym. Klasa homotopii
odwzorowania nie zależy od dostatecznie małego
wyboru otoczenia ani od metryki tubularnej.
\end{lemat}
\begin{proof}
Wybierzmy identyfikację domkniętego otoczenia
tubularnego z $D(\nu)$. Na jego wnętrzu mapę określa
rzut do ilorazu $D(\nu)/S(\nu)$; na dopełnieniu
wnętrza przyjmujemy punkt bazowy. Na brzegu tubularnym
obie formuły dają ten sam punkt, zatem z lematu o
sklejaniu odwzorowań otrzymujemy ciągłość.
Otoczenie można wybrać poza $\partial Q$, więc
odwzorowanie jest stałe na $\partial Q$ i przechodzi
na wskazany iloraz. Dwa dostatecznie małe wybory
porównujemy po zmniejszeniu do wspólnego otoczenia;
rodzina promieniowych przeskalowań włókien daje
homotopię odwzorowań ściągających. Każde dwa
otoczenia tubularne są izotopijne po zwężeniu,
co jest standardową częścią twierdzenia o otoczeniach
tubularnych. To uzasadnia niezależność klasy homotopii.
\end{proof}

Jeśli $\nu_{M/Q}$ ma uramowanie $\varphi$, otrzymujemy
złożenie
$c_{M,Q}:Q_+\to\operatorname{Th}(\nu_{M/Q})
\cong M_+\wedge S^k\to S^k$, gdzie ostatnia strzałka
zapomina punkt $M$ (na współczynniku $M_+$).
Dla $Q=S^{n+k}$ z wybranym punktem $\infty$ poza $M$
uzyskujemy mapę punktowaną $S^{n+k}\to S^k$.

\begin{center}
\begin{tikzpicture}[>=Stealth,font=\small]
\begin{scope}[xshift=0cm]
\draw[thick] (0,0) ellipse (1.50 and 0.91);
\draw[manifoldblue,thick,fill=manifoldblue!10] (0,0) ellipse (1.12 and 0.46);
\draw[manifoldred,very thick] (-0.74,-0.02) .. controls (-0.28,0.27) and (0.18,-0.24) .. (0.76,0.06);
\node[font=\footnotesize] at (0,1.20) {$S^{n+k}$};
\node[font=\footnotesize,manifoldblue] at (0,-0.70) {$U\cong D(\nu)$};
\node[font=\footnotesize,manifoldred] at (0,0.33) {$M$};
\end{scope}
\draw[->,thick] (1.65,0)--node[above,align=center,font=\footnotesize] {ściągnięcie} (4.25,0);
\begin{scope}[xshift=5.85cm]
\draw[thick,fill=manifoldblue!8] (0,0) circle (0.92);
\draw[thick,dashed] (0,0) circle (0.52);
\fill (0.88,0.14) circle (1.5pt);
\node[font=\footnotesize,anchor=west] at (1.00,0.12) {$*$};
\node[font=\footnotesize] at (0,1.20) {$\operatorname{Th}(\nu)$};
\node[font=\footnotesize] at (0,-1.18) {brzeg włókien $\mapsto *$};
\end{scope}
\end{tikzpicture}
\end{center}
\begin{center}\small Schemat idei: wnętrze rurki opisuje współrzędne normalne,
zaś jej zewnętrze zostaje utożsamione z punktem bazowym.\end{center}

\section{Twierdzenie Pontriagina dla podrozmaitości uramowanych}
\begin{definicja}[Kobordyzm uramowany w sferze]
Ustalmy $n\ge0$, $k\ge1$. Dwa zamknięte uramowane
$n$-wymiarowe podrozmaitości $M_0,M_1$ w
$S^{n+k}\setminus\{\infty\}$ są \emph{kobordantne jako podrozmaitości uramowane}
(oznaczamy $M_0\sim_{\mathrm{fr}}M_1$), jeśli istnieje
zwarta $(n+1)$-wymiarowa podrozmaitość
$W\subset(S^{n+k}\setminus\{\infty\})\times[0,1]$
prostopadła do końcowych poziomów w otoczeniach brzegów,
taka że
$\partial W=M_0\times\{0\}\sqcup M_1\times\{1\}$,
a jej wiązka normalna rzędu $k$ ma uramowanie
ograniczające się do danych uramowań końcowych.
Utożsamiamy klasy różniące się izotopią uramowaną.
Zbiór tych klas oznaczamy $\mathcal F_n^k$.
\end{definicja}

\begin{twierdzenie}[Pontriagin--Thom w ustalonym kowymiarze]
\label{thm:pt-framed-v9}
Dla $n\ge0$ i $k\ge1$ konstrukcja ściągająca
wyznacza bijekcję
\[
\operatorname{PT}:\mathcal F_n^k\xrightarrow{\cong}
[S^{n+k},S^k]_*=\pi_{n+k}(S^k).
\]
Jest zgodna ze zwiększaniem kowymiaru przez dodanie
trywialnego kierunku normalnego i z zawieszeniem
odwzorowań. W zakresie, w którym grupy homotopii
są abelowe, suma odpowiada rozłącznej sumie
uramowanych podrozmaitości umieszczonych w rozłącznych
kulach.
\end{twierdzenie}
\begin{proof}[Dowód z użyciem transwersalności i twierdzenia tubularnego]
\emph{Krok 1: od podrozmaitości do odwzorowania.}
Na mocy Lematu~\ref{lem:pt-collapse-v9} otoczenie tubularne
$M$ daje mapę $S^{n+k}\to\operatorname{Th}(\nu_{M/S^{n+k}})$.
Uramowanie utożsamia tę przestrzeń z $M_+\wedge S^k$;
po zapomnieniu parametru $M$ otrzymujemy mapę do $S^k$.
Uramowany kobordyzm $W$ ma otoczenie tubularne w
$S^{n+k}\times[0,1]$, które można wybrać produktowe
przy końcach przedziału. Jego ściągnięcie jest zatem
homotopią między mapami obu brzegów.
Wynik zależy tylko od klasy kobordyzmu.

\emph{Krok 2: od mapy do podrozmaitości.}
Niech $f:S^{n+k}\to S^k$ będzie odwzorowaniem
punktowanym. Z twierdzenia o aproksymacji gładkiej
możemy zastąpić je homotopijnym odwzorowaniem gładkim,
stałym blisko $\infty$. Wybierzmy wartość regularną
$q\ne*$ (Sard); jeśli $q$ nie należy do obrazu $f$,
wówczas $M=\varnothing$, a $f$ jest nullhomotopijne, gdyż
jego obraz mieści się w $S^k\setminus\{q\}\cong\R^k$.
W pozostałym przypadku $M=f^{-1}(q)$ jest niepustą, zwartą
$n$-podrozmaitością poza $\infty$.
Lemat~\ref{lem:pt-regular-v9} daje uramowanie normalne
przez wybór dodatnio zorientowanej bazy $T_qS^k$.
Jeśli $f_0,f_1$ są homotopijne przez homotopię punktowaną,
przybliżamy homotopię względnie do jej końców,
aby była poprzeczna do $q$. Przeciwobraz $q$
w $S^{n+k}\times[0,1]$ jest wtedy uramowanym
kobordyzmem z iloczynowym zachowaniem końców.
Tak powstaje odwzorowanie w stronę odwrotną.

\emph{Krok 3: odwrotność obu konstrukcji.}
Dla podrozmaitości z otoczeniem tubularnym mapa
ściągająca ma jako przeciwobraz punktu $q$ dokładnie
zerowy przekrój wiązki normalnej; różniczka w kierunku
normalnym odtwarza zadane uramowanie. W drugą stronę
wybierzmy małą domkniętą kulę $D$ wokół wartości
regularnej $q$, nie zawierającą punktu bazowego.
Po jej zmniejszeniu $f^{-1}(D)\to D$ jest właściwą
submersją (dziedzina i sfera są zwarte, a wszystkie
punkty blisko $q$ są regularne). Twierdzenie Ehresmanna (dla właściwej submersji,
z wersją nad brzegiem dysku) daje lokalnie trywialną
wiązkę nad $D$. Ponieważ $D$ jest ściągalny,
otrzymujemy jej trywializację $f^{-1}(D)\cong M\times D$.
Wybierając trywializację, która wzdłuż włókna centralnego
ma różniczkę normalną zgodną z $\mathrm{D}f$, uzyskujemy
identyfikację z otoczeniem tubularnym zachowującą
uramowanie po przejściu do izotopii małych dysków.
Ten ostatni krok nie wynika z samej zwartości: korzysta
z jednoznaczności otoczeń tubularnych z dokładnością
do izotopii.
Wybierzmy odwzorowanie stopnia $1$,
$S^k\to D/\partial D\cong S^k$, które wysyła
do punktu bazowego dopełnienie wnętrza $D$.
Jest ono homotopijne do identyczności (po naturalnym
utożsamieniu orientacji). Złożenie z $f$ jest więc
homotopijne do $f$, a w wybranej trywializacji
jest właśnie mapą ściągającą otoczenie tubularne.
Obie konstrukcje są wzajemnie odwrotne.

\emph{Krok 4: zawieszenie i dodawanie.}
Nowy trywialny kierunek normalny zastępuje
$D^k/S^{k-1}$ przez $D^{k+1}/S^k$;
odpowiada to zawieszeniu mapy ściągającej.
Dla dwóch rozłącznych osadzeń w kulach mapy
ściągające sumują się poprzez mapę uszczypnięcia
$S^{n+k}\to S^{n+k}\vee S^{n+k}$. To daje
zgodność z operacjami grupowymi.
\end{proof}

\begin{przyklad}[Stopień mapy jako liczba punktów]
Niech $f:S^1\to S^1$ będzie gładkim odwzorowaniem
stopnia $d$, a $q\in S^1$ jego wartością regularną.
Włókno $f^{-1}(q)$ jest skończonym zbiorem punktów.
Uramowanie normalne w każdym punkcie odpowiada znakowi
lokalnej pochodnej w zorientowanych współrzędnych.
Po zsumowaniu ze znakami otrzymujemy
$d=\sum_{x\in f^{-1}(q)}\operatorname{sgn}\D f_x$.
Tak stopień, wcześniej definiowany przez homologię,
staje się liczbą algebraiczną punktów uramowanej
zerowymiarowej podrozmaitości.
\end{przyklad}

\begin{wniosek}[Homotopia stabilna i kobordyzm uramowany]
\label{cor:pt-stable-v9}
Niech $\Omega_n^{\mathrm{fr}}:=\varinjlim_k\mathcal F_n^k$,
gdzie mapy przejścia dodają jeden trywialny kierunek
normalny. Dla każdego $n\ge0$ istnieje naturalny
izomorfizm grup
$\Omega_n^{\mathrm{fr}}\cong\pi_n^{\mathrm s}$.
\end{wniosek}
\begin{proof}
Twierdzenie~\ref{thm:pt-framed-v9} identyfikuje wszystkie
poziomy tego układu z $\pi_{n+k}(S^k)$,
a mapy przejścia z zawieszeniami. Po wzięciu
tej samej granicy prostej otrzymujemy dokładnie
definicję stabilnego pnia z poprzedniego rozdziału.
\end{proof}

\begin{przyklad}[Włókno Hopfa i pierwszy stabilny pień]
Odwzorowanie Hopfa $\eta:S^3\to S^2$ ma włókno
regularne będące okręgiem $S^1\subset S^3$, z
uramowaniem normalnym pochodzącym od $\D\eta$.
Z poprzedniego rozdziału wiemy, że $\Sigma\eta$
reprezentuje niezerowy element $\pi_1^{\mathrm s}\cong\Z/2\Z$.
Dlatego otrzymany uramowany okrąg nie jest brzegiem
uramowanego kobordyzmu, choć sam okrąg bez dodatkowej
struktury jest brzegiem dysku. Uramowanie zawiera
informację niewidoczną dla zwykłej homologii.
\end{przyklad}

\section{Związek konstrukcji z dualnością Poincarégo i klasą Eulera}
Pokażemy, dlaczego przestrzeń Thoma jest właściwym celem
odwzorowania ściągającego również wtedy, gdy normalna
wiązka nie jest trywialna. Jest to zarazem powrót do
dualności Poincarégo i teorii przecięć z wcześniejszych rozdziałów.

\begin{twierdzenie}[Ściągnięcie klasy Thoma daje dualność podrozmaitości]
\label{thm:pt-pd-v9}
Niech $Q^m$ będzie zamkniętą, spójną, zorientowaną
rozmaitością gładką, a $M^{m-k}\subset Q$ zamkniętą,
zorientowaną podrozmaitością bez brzegu i o kowymiarze
$k\ge1$. Wyposażmy $\nu=\nu_{M/Q}$ w orientację
wyznaczoną z orientacji $Q$ i $M$.
Niech $U\in\widetilde H^k(\operatorname{Th}(\nu);\Z)$
będzie całkowitą klasą Thoma, z normalizacją lokalną
zgodną z tą orientacją. Wówczas
\[
c_{M,Q}^*U=\operatorname{PD}_Q(i_*[M])
\quad\text{w }H^k(Q;\Z),
\]
gdzie $i:M\hookrightarrow Q$ jest włączeniem, a
$\operatorname{PD}_Q:H_{m-k}(Q;\Z)\to H^k(Q;\Z)$
izomorfizmem dualności Poincarégo.
Bez założeń orientowalności analogiczny wzór zachodzi
ze współczynnikami $\F_2$.
\end{twierdzenie}
\begin{proof}
Niech $\sigma:\Delta^k\to Q$ będzie gładkim
sympleksem singularnym poprzecznym do $M$ i mającym
brzeg rozłączny z $M$. Jego obraz przecina $M$ w
skończenie wielu punktach: transwersalność daje
zbiór dyskretny, a zwartość dziedziny zapewnia
skończoność. Przy każdym punkcie przecięcia możemy
wybrać współrzędne $(y,z)\in\R^{m-k}\times\R^k$
na $Q$ z $M=\{z=0\}$. Klasa Thoma w tych
współrzędnych jest generatorem względnej kohomologii
$H^k(D^k,S^{k-1};\Z)$. Wartość
$\langle c_{M,Q}^*U,\sigma\rangle$ jest więc
sumą lokalnych stopni składowej normalnej
$z\circ\sigma$ przy jej zerach; są to dokładnie
znaki poprzecznych przecięć $\sigma$ z $M$.
Wyjaśnijmy również, dlaczego otrzymujemy równość
\emph{klas całkowitych}, a nie tylko jednakowe
sparowania z cyklami wolnymi od torsji. Klasa Thoma $U$
przez ekscyzję odpowiada klasie względnej
$u\in H^k(Q,Q\setminus M;\Z)$. W łańcuchowej
konstrukcji iloczynu zewnętrznego i cap-produktu
z klasą fundamentalną lokalny rachunek ma postać
\[
[D^{m-k}\times D^k]\frown
\operatorname{pr}_{D^k}^*U_k
=[D^{m-k}\times\{0\}],
\]
gdzie $U_k$ oznacza dodatni generator
$H^k(D^k,S^{k-1};\Z)$, a klasy po obu stronach
rozumiemy w odpowiednich grupach względnych.
Istotnie, na podziale symplicjalnym zgodnym z
otoczeniem tubularnym jedynie sympleksy przecinające
$M$ wnoszą niezerowy wkład do iloczynu cap;
lokalny wkład wynosi właśnie powyższy generator.
Po zsumowaniu lokalnych fragmentów, przejściu przez
ekscyzję i zapomnieniu podpory w $Q\setminus M$
dostajemy tożsamość \emph{w homologii}
$[Q]\frown c_{M,Q}^*U=i_*[M]$.
Ponieważ $[Q]\frown(-)$ jest izomorfizmem dualności
Poincarégo, otrzymujemy żądaną równość klas
kohomologii również w obecności torsji. Współczynniki
$\F_2$ usuwają potrzebę wyboru orientacji i znaków.
\end{proof}

\begin{wniosek}[Zera przekroju i klasa Eulera]
\label{cor:euler-zero-v9}
Niech $\xi\to Q$ będzie zorientowaną rzeczywistą wiązką
wektorową rzędu $k\ge1$ nad zamkniętą, zorientowaną
rozmaitością $Q^m$. Jeżeli gładki przekrój
$s:Q\to\xi$ jest poprzeczny do przekroju zerowego,
to $Z(s)=\{x:s(x)=0\}$ jest zamkniętą, zorientowaną
podrozmaitością wymiaru $m-k$ i
$\operatorname{PD}_Q([Z(s)])=e(\xi)$ w $H^k(Q;\Z)$.
\end{wniosek}
\begin{proof}
Transwersalność przekroju do przekroju zerowego
daje podrozmaitość wymiaru $m-k$, a pochodna przekroju
w kierunkach normalnych identyfikuje
$\nu_{Z(s)/Q}$ z $\xi|_{Z(s)}$ i przenosi orientację.
Definicja klasy Eulera jako przeciwobrazu klasy Thoma
przez przekrój zerowy i homotopia $t\mapsto t s$
utożsamiają ją z lokalną klasą zer przekroju.
Twierdzenie~\ref{thm:pt-pd-v9} identyfikuje tę klasę
z dualną Poincarégo do $[Z(s)]$. Jeśli $k>m$,
transwersalny przekrój nie ma zer; wówczas obie strony
są zerowe, ponieważ $H^k(Q;\Z)=0$.
\end{proof}

\begin{przyklad}[Dlaczego znak ma znaczenie]
Dla zorientowanej wiązki liniowej nad okręgiem
każdy przekrój poprzeczny do zera ma zera o
lokalnych indeksach $+1$ lub $-1$, których suma
jest zerowa. Odpowiada to $H^1(S^1;\Z)\cong\Z$,
ale klasa Eulera zorientowanej rzeczywistej wiązki
rzędu $1$ jest zerowa: orientowana wiązka liniowa
jest trywialna. Bez orientacji można natomiast otrzymać
wiązki Möbiusa i niezerową klasę $w_1$.
\end{przyklad}

\section{Kobordyzmy nieorientowane i zorientowane}
\begin{definicja}[Kobordyzm gładki]
Niech $M_0,M_1$ będą zamkniętymi gładkimi
$n$-rozmaitościami. \emph{Kobordyzmem (nieorientowanym)}
między nimi nazywamy zwartą gładką
$(n+1)$-rozmaitość $W$ wraz z dyfeomorfizmem
$\partial W\cong M_0\sqcup M_1$.
W kategorii zorientowanej wymagamy orientacji wszystkich
obiektów i izomorfizmu zorientowanych brzegów
$\partial W\cong (-M_0)\sqcup M_1$, gdzie
$-M_0$ ma odwróconą orientację, a $\partial W$ ma
orientację indukowaną konwencją normalnej zewnętrznej.
Zbiory klas oznaczamy odpowiednio
$\Omega_n^O$ i $\Omega_n^{SO}$. Dodawanie klas
jest zadane sumą rozłączną, a mnożenie produktem
kartezjańskim rozmaitości.
\end{definicja}

\begin{stwierdzenie}[Struktury algebraiczne]
\label{prop:bordism-group-v9}
Dla każdego $n\ge0$ zbiory $\Omega_n^O$ i
$\Omega_n^{SO}$ są grupami abelowymi. W grupie
nieorientowanej zachodzi $2[M]=0$, a w grupie
zorientowanej $[-M]=-[M]$. Sumy
$\Omega_*^O:=\bigoplus_{n\ge0}\Omega_n^O$
i $\Omega_*^{SO}:=\bigoplus_{n\ge0}\Omega_n^{SO}$
są pierścieniami stopniowanymi, przy czym
w przypadku zorientowanym mnożenie jest
przemienne z regułą znaku $(-1)^{mn}$ w stopniach
$m,n$.
\end{stwierdzenie}
\begin{proof}
Sumowanie rozłączne jest łączne i przemienne z
dokładnością do dyfeomorfizmu; pusty zbiór jest
zerem. Walec $M\times[0,1]$ ma dwa egzemplarze
$M$ jako brzeg, co daje $[M]+[M]=0$ bez orientacji.
Z orientacją brzeg walca ma postać $(-M)\sqcup M$,
więc $[-M]+[M]=0$. Jeśli $M\sim M'$ przez $W$,
to dla zamkniętego $N$ produkt $W\times N$ jest
kobordyzmem między $M\times N$ i $M'\times N$.
Zatem mnożenie nie zależy od reprezentantów.
Zamiana czynników $M^m\times N^n\to N^n\times M^m$
zmienia orientację o $(-1)^{mn}$, co wyjaśnia
regułę przemienności. Pozostałe aksjomaty wynikają
z łączności i rozdzielności produktu względem
sumy rozłącznej.
\end{proof}

\begin{przyklad}[Pierwsze grupy kobordyzmu]
\label{ex:bordism-low-v9}
Wymiar zero daje $\Omega_0^O\cong\Z/2\Z$
(parzystość liczby punktów) oraz
$\Omega_0^{SO}\cong\Z$ (liczba punktów zorientowanych,
ze znakami). Każda zwarta $1$-rozmaitość ma parzystą
liczbę punktów brzegowych, bo jej składowe z brzegiem są
odcinkami. I odwrotnie, parę punktów można połączyć
odcinkiem, zorientowanym od punktu ujemnego do dodatniego
w przypadku zorientowanym. To dowodzi obu izomorfizmów. Każda zwarta rozmaitość jednowymiarowa
bez brzegu jest sumą okręgów, a okrąg jest brzegiem
dysku, więc $\Omega_1^O=\Omega_1^{SO}=0$.
Każda zamknięta zorientowana powierzchnia jest brzegiem
trójwymiarowego ciała z uchwytami, dlatego
$\Omega_2^{SO}=0$. Z klasyfikacji zamkniętych
powierzchni otrzymamy dalej
$\Omega_2^O\cong\Z/2\Z$, generowaną przez
$\R P^2$. W dwóch ostatnich wnioskach używamy
klasyfikacji powierzchni jako wyniku znanego z
topologii rozmaitości, a nie twierdzenia dowiedzionego
tutaj od początku.
\end{przyklad}

\begin{przyklad}[Torus i butelka Kleina: różne topologie, ten sam kobordyzm]
\label{prz:bordism-torus-klein-24}
Zbudujemy dwie zwarte gładkie trójrozmaitości, których brzegami są odpowiednio
$T^2$ i butelka Kleina $K$. Niech $D^2=\{z\in\C:|z|\le1\}$ oraz
$r:D^2\to D^2$, $r(z)=\overline z$. Określamy
\[
 V_+=D^2\times S^1,\qquad
 V_-=(D^2\times[0,1])/((z,1)\sim(r(z),0)).
\]
Drugą przestrzeń nazywamy skręconą wiązką dysków nad okręgiem.
Sklejenie końców walca za pomocą dyfeomorfizmu daje gładką
rozmaitość z brzegiem. Na podzbiorze $\partial D^2\times[0,1]$
relacja daje
\[
 \partial V_+=S^1\times S^1=T^2,\qquad
 \partial V_-=(S^1\times[0,1])/((z,1)\sim(\overline z,0))=K.
\]
Ponieważ oba obiekty są brzegami, mamy $[T^2]=[K]=0$
w $\Omega_2^O$, bez odwoływania się do klasyfikacji powierzchni.
Jednak ich topologie się różnią: jeśli $a$ oznacza obieg włókna,
a $b$ obieg podstawy, to z twierdzenia van Kampena
\[
 \pi_1(T^2)=\langle a,b\mid bab^{-1}=a\rangle,\qquad
 \pi_1(K)=\langle a,b\mid bab^{-1}=a^{-1}\rangle.
\]
Po abelianizacji pierwsza prezentacja daje $\Z^2$, druga zaś
$\Z\oplus\Z/2\Z$ (relacja $2[a]=0$). Dokładniej,
$H_1(T^2;\Z)\cong\Z^2$ oraz
$H_1(K;\Z)\cong\Z\oplus\Z/2\Z$. Nie istnieje zatem homeomorfizm
$T^2\cong K$. Mimo to $\chi(T^2)=\chi(K)=0$ i ich klasy
nieorientowanego kobordyzmu są identyczne. Pokazuje to, że
kobordyzm nie zachowuje całkowitych grup homologii.
Konstrukcję tę można porównać z rachunkiem uchwytów w
\href{https://github.com/dariusz-piekarz/mathematical_scipts/blob/biblioteka-matematyczna-iteracja-6/chapters/pl/differential_geometry_and_topology/19-kobordyzmy-i-rachunek-uchwytow.tex}{rozdziale o kobordyzmach i uchwytach}:
różne przyklejenia jednego $1$-uchwytu dają ten sam kontrast.
\end{przyklad}

\section{Klasy charakterystyczne jako przeszkody kobordyzmu}
\begin{definicja}[Liczba Stiefela--Whitneya]
Niech $M^n$ będzie zamkniętą gładką $n$-rozmaitością,
a $w_j(TM)\in H^j(M;\F_2)$ klasami Stiefela--Whitneya.
Dla nieujemnych indeksów $i_1,\ldots,i_s$ o sumie
$i_1+\cdots+i_s=n$ liczbę
\[
\bigl\langle w_{i_1}(TM)\smile\cdots\smile
w_{i_s}(TM),[M]_2\bigr\rangle\in\F_2
\]
nazywamy \emph{liczbą Stiefela--Whitneya}.
Tutaj $[M]_2\in H_n(M;\F_2)$ oznacza klasę fundamentalną;
nie wymaga ona orientowalności $M$.
\end{definicja}

\begin{twierdzenie}[Zanik liczb charakterystycznych na brzegu]
\label{thm:sw-bdy-v9}
Jeśli zamknięta $n$-rozmaitość $M$ jest dyfeomorficzna
z brzegiem zwartej $(n+1)$-rozmaitości $W$, to
wszystkie liczby Stiefela--Whitneya $M$ są zerowe.
Jeżeli $M$ i $W$ są zorientowane, analogicznie
zanikają liczby Pontriagina złożone z klas
$p_j(TM)\in H^{4j}(M;\Z)$ o łącznym stopniu $n$.
\end{twierdzenie}
\begin{proof}
Niech $i:M\hookrightarrow W$ oznacza włączenie brzegu.
Z twierdzenia o kołnierzu otrzymujemy
$TW|_M\cong TM\oplus\underline\R$.
Stabilność klas charakterystycznych daje
$i^*w_j(TW)=w_j(TM)$ oraz analogicznie dla $p_j$.
Jeżeli $P$ jest dowolnym jednomianem charakterystycznym
stopnia $n$, to
$\langle P(TM),[M]\rangle=
\langle i^*P(TW),[M]\rangle=
\langle P(TW),i_*[M]\rangle$.
Ale $i_*[M]=0$ w $H_n(W;\F_2)$:
wynika to z długiego ciągu homologii pary $(W,M)$,
ponieważ $[M]$ jest obrazem względnej klasy fundamentalnej
$[W,M]\in H_{n+1}(W,M;\F_2)$ przez homomorfizm brzegowy.
W przypadku zorientowanym ten sam argument prowadzimy
ze współczynnikami $\Z$, używając zgodności
orientacji brzegu. Każde rozważane sparowanie wynosi zero.
\end{proof}

\begin{przyklad}[Dlaczego $\R P^2$ nie jest brzegiem]
Niech $t\in H^1(\R P^2;\F_2)$ będzie generatorem.
W punkcie odpowiadającym prostej $\ell\subset\R^{n+1}$
przestrzeń styczna $T_{[\ell]}\R P^n$ utożsamia się
z $\operatorname{Hom}(\ell,\ell^\perp)$.
Dla skończenie wymiarowych przestrzeni liniowych
$\operatorname{Hom}(\ell,\ell^\perp)\cong\ell^*\otimes\ell^\perp$.
Wyprowadzenie tego izomorfizmu oraz własność uniwersalna
iloczynu tensorowego znajdują się w
\href{https://github.com/dariusz-piekarz/mathematical_scipts/blob/biblioteka-matematyczna-iteracja-6/chapters/pl/differential_geometry_and_topology/03-algebra-tensorowa-i-zewnetrzna.tex}{rozdziale \emph{Algebra tensorowa i zewnętrzna}}
skryptu z geometrii różniczkowej
(etykiety źródłowe \texttt{ch:algebra-tensorowa},
\texttt{sec:iloczyn-tensorowy}, \texttt{thm:tensor-hom}).
Z rozkładu
$\R^{n+1}=\ell\oplus\ell^\perp$ wynika izomorfizm
wiązek
$\underline\R\oplus T\R P^n
\cong\operatorname{Hom}(\gamma^1,\underline\R^{n+1})
\cong(n+1)(\gamma^1)^*$.
Metryka euklidesowa utożsamia $(\gamma^1)^*$ z
$\gamma^1$. Stąd
$T\R P^n\oplus\underline\R\cong(n+1)\gamma^1$,
a więc $w(T\R P^2)=(1+t)^3=1+t+t^2$.
Stąd $\langle w_1^2,[\R P^2]_2\rangle=1$ i
$\langle w_2,[\R P^2]_2\rangle=1$, ponieważ
$t^2$ generuje $H^2(\R P^2;\F_2)$.
Na mocy Twierdzenia~\ref{thm:sw-bdy-v9}
$\R P^2$ nie jest brzegiem zwartej $3$-rozmaitości.
Warto przy tym pamiętać, że $\R P^2$ nie jest
orientowalna: argument dotyczy \emph{nieorientowanego}
kobordyzmu.
\end{przyklad}

\begin{przyklad}[Niezerowa klasa kobordyzmu produktu $\R P^2\times\R P^2$]
\label{prz:bordism-rp2-product}
Niech $X=\R P^2\times\R P^2$ oraz
$a=\operatorname{pr}_1^*t$, $b=\operatorname{pr}_2^*t$, gdzie
$t\in H^1(\R P^2;\F_2)$ jest generatorem. Wzór Künnetha daje
\[
 H^*(X;\F_2)=\F_2[a,b]/(a^3,b^3),\qquad
 \deg a=\deg b=1.
\]
Z obliczenia $w(T\R P^2)=(1+t)^3=1+t+t^2$, naturalności
oraz wzoru Whitney'a otrzymujemy
\[
 w(TX)=(1+a+a^2)(1+b+b^2),\quad
 w_2=a^2+ab+b^2,\quad w_4=a^2b^2.
\]
W pierścieniu charakterystyki dwa zachodzi
$w_2^2=a^4+a^2b^2+b^4=a^2b^2$.
Ponieważ $a^2b^2$ generuje $H^4(X;\F_2)$, obie liczby
\[
 \langle w_2(TX)^2,[X]_2\rangle=1,\qquad
 \langle w_4(TX),[X]_2\rangle=1
\]
są niezerowe. Twierdzenie~\ref{thm:sw-bdy-v9} wyklucza więc
istnienie zwartej pięciorozmaitości z brzegiem $X$.
Jednocześnie $2[X]=0$ w $\Omega_4^O$ na mocy
Stwierdzenia~\ref{prop:bordism-group-v9}; zatem $[X]$ ma dokładnie
rząd dwa. Nie wyznacza to jeszcze całej grupy $\Omega_4^O$.
\end{przyklad}

\begin{przyklad}[Ta sama parzystość Eulera, różne klasy kobordyzmu]
\label{prz:bordism-rp4-versus-rp2-square}
Niech $X=\R P^2\times\R P^2$ i $Y=\R P^4$.
Obie przestrzenie mają charakterystykę Eulera równą jeden:
$\chi(X)=\chi(\R P^2)^2=1=\chi(Y)$.
Mimo to nie są kobordantne. Istotnie, dla generatora
$t\in H^1(Y;\F_2)$ mamy
\[
 w(TY)=(1+t)^5=1+t+t^4\pmod {t^5},
\]
gdyż $\binom52$ i $\binom53$ są parzyste, natomiast
$\binom51$ i $\binom54$ nieparzyste. W szczególności
$w_2(TY)=0$, $w_4(TY)=t^4$. Porównując liczby
charakterystyczne otrzymujemy
\[
 \bigl(\langle w_2^2,[X]_2\rangle,\langle w_4,[X]_2\rangle\bigr)
  =(1,1),\quad
 \bigl(\langle w_2^2,[Y]_2\rangle,\langle w_4,[Y]_2\rangle\bigr)
  =(0,1).
\]
Gdyby $X$ i $Y$ były kobordantne, suma ich klas byłaby
zerem, więc każda liczba Stiefela--Whitneya byłaby na nich
jednakowa (twierdzenie o zaniku na brzegu). Otrzymujemy
sprzeczność. Co więcej, wyznaczone wektory $(1,1)$ i $(0,1)$
są liniowo niezależne w $(\F_2)^2$, a klasy $[X]$ i $[Y]$
mają rząd dwa. Podgrupa przez nie generowana jest zatem
izomorficzna z $(\Z/2\Z)^2$. To dolne ograniczenie rozmiaru
$\Omega_4^O$, a nie pełne obliczenie tej grupy.
\end{przyklad}

\begin{figure}[!htbp]
\centering
\begin{tikzpicture}[>=Stealth,font=\small]
 \node[draw,rounded corners,fill=manifoldblue!8,
       minimum width=4cm,minimum height=1.05cm,
       align=center] (a) at (0,0)
       {$w(T\R P^2)=1+t+t^2$};
 \node[draw,rounded corners,fill=manifoldgreen!8,
       minimum width=4.5cm,minimum height=1.05cm,
       align=center] (b) at (6,0)
       {$w(TX)=(1+a+a^2)(1+b+b^2)$};
 \draw[->,thick] (a)--(b)
       node[midway,above=7pt,font=\scriptsize] {Whitney};
 \node[draw,rounded corners,fill=manifoldgold!10,
       minimum width=4.5cm,minimum height=1.05cm,
       align=center] (c) at (6,-2)
       {$w_2^2=w_4=a^2b^2$\\
        $\langle a^2b^2,[X]_2\rangle=1$};
 \node[draw,rounded corners,fill=manifoldred!8,
       minimum width=4cm,minimum height=1.05cm,
       align=center] (d) at (0,-2)
       {$[X]\ne0$ w $\Omega_4^O$\\$2[X]=0$};
 \draw[->,thick] (b)--(c);
 \draw[->,thick] (c)--(d);
\end{tikzpicture}
\caption{Liczby charakterystyczne produktu
$\R P^2\times\R P^2$ wykrywają niezerowy element
rzędu dwa w nieorientowanym kobordyzmie.}
\label{fig:bordism-rp2-product}
\end{figure}

\begin{twierdzenie}[Obliczenie $\Omega_2^O$]
\label{thm:omega2-v9}
Istnieje izomorfizm
$\Omega_2^O\cong\Z/2\Z$ zadany przez
$[M]\mapsto\chi(M)\bmod2$; generatorem jest
klasa $\R P^2$.
\end{twierdzenie}
\begin{proof}
Dla zamkniętej powierzchni $M$ równość
$\langle w_2(TM),[M]_2\rangle=\chi(M)\bmod2$
wynika z twierdzenia Poincarégo--Hopfa modulo $2$.
Na mocy Twierdzenia~\ref{thm:sw-bdy-v9} odwzorowanie
jest więc niezmiennikiem kobordyzmu.
Każda zamknięta spójna powierzchnia jest orientowalną
powierzchnią $\Sigma_g$ albo sumą spójną $h$ egzemplarzy
$\R P^2$. Powierzchnie $\Sigma_g$ są brzegami ciał
z uchwytami. Z kolei weźmy
$(M\sqcup N)\times[0,1]$ i dołączmy do jego górnej
części trójwymiarowy uchwyt $D^1\times D^2$
wzdłuż dwóch rozłącznych dysków $D^2$,
po jednym na każdej powierzchni. Nowa górna
powierzchnia jest sumą spójną $M\#N$,
dolna pozostaje $M\sqcup N$. W ten sposób
otrzymujemy kobordyzm dowodzący
$[M\#N]=[M]+[N]$ w $\Omega_2^O$.
Wynika stąd, że każda klasa jest wielokrotnością
$[\R P^2]$. Równanie $2[\R P^2]=0$ pochodzi
ze Stwierdzenia~\ref{prop:bordism-group-v9}, a liczba
$\chi(\R P^2)\bmod2=1$ pokazuje, że generator
jest niezerowy. Dostajemy żądany izomorfizm.
\end{proof}

\begin{uwaga}[Zakres twierdzenia Thoma o liczbach charakterystycznych]
W dalszej teorii dowodzi się głębszego twierdzenia:
\emph{dwie zamknięte gładkie $n$-rozmaitości są
nieorientowanie kobordantne wtedy i tylko wtedy, gdy
wszystkie odpowiadające im liczby Stiefela--Whitneya
są równe}. Kierunek konieczny udowodniliśmy powyżej;
kierunek dostateczny opiera się na obliczeniu pierścienia
$\Omega_*^O$ metodami Thoma i nie jest tu dowiedziony.
\end{uwaga}

\section{Od uramowania do widm \texorpdfstring{$MO$ i $MSO$}{MO i MSO}}
\begin{definicja}[Przestrzenie Grassmanna i wiązki tautologiczne]
Dla $k\ge1$, $N\ge0$ oznaczmy przez
$\mathrm{Gr}_k(\R^{k+N})$ przestrzeń $k$-wymiarowych
podprzestrzeni liniowych $\R^{k+N}$. Jej punktowi $V$
odpowiada włókno $V$ rzeczywistej wiązki tautologicznej
$\gamma_{k,N}$. Granica włączeń
$\mathrm{Gr}_k(\R^{k+N})\hookrightarrow
\mathrm{Gr}_k(\R^{k+N+1})$ daje model $BO(k)$
z uniwersalną wiązką $\gamma_k$. Zastąpienie
$\mathrm{Gr}_k$ przestrzenią zorientowanych
$k$-płaszczyzn daje $BSO(k)$ i $\gamma_k^{SO}$.
Przestrzenie Thoma tych wiązek oznaczamy
$MO(k):=\operatorname{Th}(\gamma_k)$ oraz
$MSO(k):=\operatorname{Th}(\gamma_k^{SO})$.
\end{definicja}

\begin{stwierdzenie}[Klasyfikacja wiązki normalnej]
\label{prop:normal-class-v9}
Jeśli $M^n\subset\R^{n+k}$ jest zamkniętą
podrozmaitością, to jej wiązka normalna
$\nu_M$ jest klasyfikowana przez odwzorowanie
$c_\nu:M\to BO(k)$; w szczególności
$\nu_M\cong c_\nu^*\gamma_k$.
Jeśli $M$ jest zorientowana, to $\nu_M$ jest
orientowalna i można wybrać podniesienie mapy
klasyfikującej do $BSO(k)$, zgodne z orientacją
$TM\oplus\nu_M\cong\underline\R^{n+k}$.
\end{stwierdzenie}
\begin{proof}
Wybierzmy metrykę euklidesową. Włókno normalne
$\nu_{M,x}$ jest $k$-wymiarową podprzestrzenią
$\R^{n+k}$, więc funkcja $x\mapsto\nu_{M,x}$
jest ciągłą (a nawet gładką) mapą do
$\mathrm{Gr}_k(\R^{n+k})$. Jej cofnięciem
$\gamma_{k,n}$ jest z definicji właśnie $\nu_M$.
Wybrana orientacja $M$ i standardowa orientacja
$\R^{n+k}$ wyznaczają orientację każdego włókna
normalnego przez zasadę sumy bezpośredniej.
Stąd otrzymujemy mapę do zorientowanej
przestrzeni Grassmanna.
\end{proof}

\begin{twierdzenie}[Pontriagin--Thom dla kobordyzmu]
\label{thm:pt-MO-v9}
Dla każdego $n\ge0$ istnieją naturalne izomorfizmy
\[
\Omega_n^O\cong\pi_n(MO):=
\varinjlim_{k\to\infty}\pi_{n+k}(MO(k)),
\qquad
\Omega_n^{SO}\cong\pi_n(MSO):=
\varinjlim_{k\to\infty}\pi_{n+k}(MSO(k)).
\]
Mapy przejścia po prawej pochodzą z odwzorowań
$\Sigma MO(k)\to MO(k+1)$ (odpowiednio
$\Sigma MSO(k)\to MSO(k+1)$), wyznaczonych
przez dodanie trywialnej prostej do wiązki.
Po prawej występują grupy homotopii \emph{widm},
a nie grupy homotopii pojedynczej przestrzeni $BO$
czy $BSO$.
\end{twierdzenie}
\begin{proof}
\emph{Od rozmaitości do mapy.}
Osadźmy zamkniętą gładką $n$-rozmaitość
$M$ w $\R^{n+k}$ dla dostatecznie dużego $k$
(twierdzenie Whitneya). Osadzenie ma otoczenie
tubularne oraz wiązkę normalną $\nu_M$.
Po uzwartowieniu jednopunktowym
$\R^{n+k}\cup\{\infty\}=S^{n+k}$ otrzymujemy mapę ściągającą
$S^{n+k}\to\operatorname{Th}(\nu_M)$.
Wybór mapy klasyfikującej
$c_\nu:M\to BO(k)$ daje odwzorowanie przestrzeni
Thoma $\operatorname{Th}(\nu_M)\to MO(k)$,
a złożenie definiuje klasę w $\pi_{n+k}(MO(k))$.
Dla rozmaitości zorientowanej wybieramy orientację normalną
wyznaczoną przez standardową orientację przestrzeni
$\R^{n+k}$ i orientację $M$; wówczas mapą docelową jest
$MSO(k)$. Nie chodzi tu o trywializację wiązki normalnej.
Zmiana osadzenia w dostatecznie dużym wymiarze nie
zmienia stabilnej klasy. Dokładniej, jeśli
$i_0,i_1:M\hookrightarrow\R^{n+k}$ są dwoma osadzeniami,
po przejściu do $\R^{2(n+k)+1}$ można zestawić
$i_0$ i $i_1$ w rozłącznych ortogonalnych składnikach
i połączyć przez rodzinę osadzeń
$x\mapsto((1-t)i_0(x),t i_1(x),t)$; dla $t\in(0,1)$
różnowartościowość i iniektywność różniczki wynikają
z odpowiednich własności $i_0$ i $i_1$.
Po zaokrągleniu końców rodzina wyznacza izotopię
w wymiarze stabilnym. Zachowanie orientacji
normalnej wynika z ciągłości orientacji w rodzinie.
Identyfikacja nadmiarowych kierunków normalnych
jest zgodna z mapami strukturalnymi widma.
Jeśli $M_0$ i $M_1$ są kobordantne przez $W$,
możemy po zwiększeniu wymiaru osadzić $W$ w
$\R^{n+k}\times[0,1]$, produktowo przy brzegu.
Ściągnięcie otoczenia tubularnego $W$ daje homotopię
między klasami z $M_0$ i $M_1$. Tym samym
otrzymujemy dobrze określony homomorfizm z
odpowiedniej grupy kobordyzmu.

\emph{Od mapy do rozmaitości.}
Klasa $\pi_{n+k}(MO(k))$ jest reprezentowana
przez mapę z \emph{kompaktowej} sfery. W modelu
CW przestrzeni $MO(k)$ można, po homotopii
komórkowej, sprowadzić tę mapę do przestrzeni
Thoma pewnej skończenie wymiarowej wiązki
$\gamma_{k,N}$ nad $\mathrm{Gr}_k(\R^{k+N})$.
Nie wolno traktować całej przestrzeni Thoma jako
rozmaitości gładkiej w punkcie bazowym: jest ona
ilorazem, a więc ten krok wymaga lokalizacji.
Wybieramy otwarte otoczenie przekroju zerowego
we wnętrzu $D(\gamma_{k,N})$; na jego przeciwobrazie
mapa do przestrzeni Thoma przyjmuje wartości
w gładkiej całkowitej przestrzeni wiązki.
Stosujemy względne twierdzenie o aproksymacji
i transwersalności na zwartym obszarze zawierającym
przeciwobraz przekroju zerowego, pozostawiając
odwzorowanie niezmienione w sąsiedztwie punktu bazowego.
Po tej operacji wybieramy reprezentanta poprzecznego
do przekroju zerowego. Przeciwobraz sekcji zerowej
jest zamkniętą gładką
$n$-podrozmaitością $M\subset S^{n+k}$, a jej
wiązka normalna jest izomorficzna z cofnięciem
$\gamma_{k,N}$. W wersji zorientowanej analogiczna
konstrukcja zachowuje orientację normalną.
Homotopię map można uczynić poprzeczną do sekcji
zerowej względnie do końców; jej przeciwobraz jest
kobordyzmem otrzymanych rozmaitości.

\emph{Sprawdzenie odwrotności i stabilizacji.}
Dla map ściągających przeciwobraz sekcji zerowej
to początkowa rozmaitość wraz z daną wiązką normalną.
Odwrotnie, w otoczeniu regularnego przeciwobrazu
przybliżona mapa do przestrzeni Thoma daje współrzędne
normalne; ściągnięcie otoczenia tubularnego odtwarza
jej klasę homotopii, jak w dowodzie
Twierdzenia~\ref{thm:pt-framed-v9}.
Dodanie trywialnego kierunku normalnego jest zgodne
z $\Sigma MO(k)\to MO(k+1)$. Przechodząc do granic
prostych, uzyskujemy obustronne izomorfizmy.
\end{proof}

\begin{uwaga}[Co zostało użyte jako narzędzie pomocnicze?]
Dowód zawiera konstrukcję obu odwzorowań i argumenty,
że są wzajemnie odwrotne. Wymaga jednak standardowych
wyników geometrii różniczkowej: osadzania Whitneya,
otoczeń tubularnych, aproksymacji gładkiej,
transwersalności względnej i stabilizacji osadzeń.
Ich pełnych dowodów nie powtarzamy; znajdują się
w osobnej części książki o rozmaitościach.
Porównajmy zwłaszcza
\href{https://github.com/dariusz-piekarz/mathematical_scipts/blob/biblioteka-matematyczna-iteracja-6/chapters/pl/differential_geometry_and_topology/19-kobordyzmy-i-rachunek-uchwytow.tex}{rozdział \emph{Kobordyzmy i rachunek uchwytów}}
(etykieta źródłowa \texttt{ch:rachunek-uchwytow}). Podobnie
techniczny argument faktoryzacji przez skończony
etap Grassmanna jest faktem o filtracji CW.
Nie jest to dowód obliczenia całych pierścieni
$\Omega_*^O$ i $\Omega_*^{SO}$.
\end{uwaga}


\section{Uzupełnienia audytu: zakres transwersalności i konstrukcje stabilne}
\begin{uwaga}[Wartość regularna poza obrazem]
Dla mapy $f:S^{n+k}\to S^k$ wartość regularna może nie należeć
do obrazu. Jeżeli $q\notin f(S^{n+k})$, przeciwobraz $f^{-1}(q)$
jest pustą podrozmaitością, a mapa $f$ jest homotopijna ze stałą,
ponieważ $S^k\setminus\{q\}$ jest ściągalna. Rozróżnienie
to jest niezbędne, aby konstrukcja odwrotna Pontriagina--Thoma
obejmowała również klasę zerową.
\end{uwaga}

\begin{lemat}[Podstawowy przypadek aproksymacji dla przestrzeni Thoma]
\label{lem:thom-transversal-v9-audyt}
Niech $\xi\to B$ będzie gładką wiązką rzędu $k$ nad zwartą
gładką rozmaitością $B$, a $h:S^N\to\operatorname{Th}(\xi)$
ciągłym odwzorowaniem punktowanym. Wtedy $h$ jest
homotopijne przez homotopię zachowującą punkt bazowy,
z odwzorowaniem, którego przeciwobraz $B\subset\operatorname{Th}(\xi)$
przekroju zerowego jest gładką podrozmaitością
kowymiaru $k$ z normalną wiązką izomorficzną do
cofnięcia $\xi$.
\end{lemat}
\begin{proof}
Wybierzmy mniejsze otwarte otoczenie $V$ przekroju zerowego
w $D(\xi)$, którego domknięcie nie przecina $S(\xi)$.
Jego obraz w ilorazie jest gładką otwartą podrozmaitością
$\operatorname{Th}(\xi)$, ponieważ iloraz nie utożsamia
żadnych punktów $V$. Zbiór $h^{-1}(B)$ jest domknięty,
więc zwarty. Ma otoczenie $U\subset S^N$, którego
obraz mieści się w $V$. Twierdzenie o względnej gładkiej
aproksymacji pozwala wygładzić $h|_U$ w mniejszym
otoczeniu $h^{-1}(B)$, zaś twierdzenie o transwersalności
pozwala wybrać małe zaburzenie poprzeczne do przekroju
zerowego. Zaburzenie może mieć podporę zwartą w $U$:
mapa pozostaje zatem niezmieniona poza $U$, w tym w punkcie bazowym (po wyborze $U$ rozłącznego
z jego przeciwobrazem). Twierdzenie o przeciwobrazie
transwersalnym daje podrozmaitość kowymiaru $k$.
Różniczka poprzecznego odwzorowania identyfikuje jej
normalną wiązkę z cofnięciem normalnej przekroju
zerowego, a ta ostatnia jest kanonicznie izomorficzna
z $\xi$. Cały argument używa zwykłej transwersalności
\emph{tylko na gładkim obszarze przestrzeni Thoma},
nie na jej punkcie bazowym.
\end{proof}

\begin{uwaga}[Podstawa i wybór grupy strukturalnej]
Orientacja $M$ nie daje uramowania $\nu_M$;
daje jedynie redukcję grupy strukturalnej normalnej
wiązki z $O(k)$ do $SO(k)$ przy wyborze orientacji otoczenia.
Dlatego klasyfikacja przez $MSO$ nie jest szczególnym
przypadkiem klasyfikacji uramowanej bez zapomnienia
dodatkowych danych. Z kolei dla ogólnej przestrzeni
klasyfikującej $BO(k)$ wymaga się numerycznych wiązek;
na zwartych rozmaitościach użytych tutaj warunek ten
jest spełniony.
\end{uwaga}

\section{Po co nam widma Thoma?}
Zbierzmy konstrukcje w jednym diagramie pojęciowym:
\[
\begin{gathered}
\text{uramowanie normalne} \quad\longleftrightarrow\quad
\pi_*^{\mathrm s}\cong\Omega_*^{\mathrm{fr}},\\[3pt]
\text{wiązka normalna bez orientacji}\quad\longleftrightarrow\quad
\pi_*(MO)\cong\Omega_*^O,\\[3pt]
\text{zorientowana wiązka normalna}\quad\longleftrightarrow\quad
\pi_*(MSO)\cong\Omega_*^{SO}.
\end{gathered}
\]
Widzimy, że \emph{różne rodzaje struktur normalnych}
prowadzą do różnych teorii kobordyzmu i różnych widm.
Tę samą ideę można zastosować do struktur zespolonych;
powstaje widmo $MU$ i teoria kobordyzmu zespolonego
$\Omega_*^U\cong\pi_*(MU)$. W tym przypadku chodzi
o stabilną strukturę zespoloną na wiązce normalnej;
nie wymagamy, aby sama rozmaitość miała strukturę
rozmaitości zespolonej. Podobnie struktury spinowe
prowadzą do widma $MSpin$, ale wymagają wyjaśnienia
odpowiednich przeszkód kohomologicznych.

\begin{uwaga}[Późniejsze połączenie z geometrią algebraiczną]
Kobordyzm zespolony topologiczny $MU$ jest szczególnie
ważny dla topologii przestrzeni zespolonych. Nie należy
go jednak \emph{utożsamiać bez dodatkowej konstrukcji}
z algebraiczną teorią kobordyzmu na rozmaitościach
algebraicznych (np. teorią Levine'a--Morela).
Porównanie wymaga map realizacji i hipotez dotyczących
ciała bazowego, gładkości oraz użytej teorii orientowanej.
\end{uwaga}

\section{Podsumowanie i dalsza droga}
Otrzymaliśmy rzeczywiste przejście między językiem
homotopii i geometrią podrozmaitości: wartości regularne
tworzą podrozmaitości uramowane, a otoczenia tubularne
zamieniają je z powrotem w odwzorowania do sfer.
Przestrzenie Thoma pozwalają ten mechanizm uogólnić
na kobordyzmy bez ramowania. Wersje $MO$ i $MSO$
dają geometryczną interpretację grup homotopii widm,
a klasy charakterystyczne wykrywają przeszkody do
zerowania klas kobordyzmu.

W dalszych etapach naturalne są: szczegółowe obliczenie
pierścienia $\Omega_*^O$, teoria klas Stiefela--Whitneya
jako pełnych niezmienników kobordyzmu, orientowalność
zespolona widm $MU$ i współczynniki formalnych praw
grupowych. Kolejne dwa rozdziały poświęcamy homologii Morse'a oraz
twierdzeniu Lefschetza o punktach stałych.

% Literatura do rozdziału: R. Cohen, Bundles, Manifolds and Homotopy,
% rozdział o Pontriaginie--Thomie; R. E. Stong, Notes on Cobordism Theory.

% Iteracja X, rozdział XVIII. Wymaga preambuły wersji IX.